% Source: MQ12a_v1.html
\documentclass[11pt]{article}
\usepackage[margin=1in]{geometry}
\usepackage{amsmath}
\usepackage{amssymb}
\usepackage{siunitx}
\usepackage{enumitem}

\title{MQ12a: Fluid Mechanics}
\author{}
\date{}

\begin{document}
\maketitle

\section*{MQ12a: Fluid Mechanics}

\begin{enumerate}[leftmargin=*,label=\textbf{Question \arabic*.}]
\item \hfill \textit{1 pts}

Suppose the density of object A is 0.5 times that of object B, whereas the volume of object A is 2.7 times that of object B. If the mass of A is 1.7 kg, calculate the mass of object B in kilograms.



\emph{m }= \rule{2cm}{0.4pt} kg

\item \hfill \textit{1 pts}

A hollow cylindrical pipe is made from material 3 listed in the table below. The pipe is 3.3 m long, has an outside diameter of 3.6 cm, and an inside diameter of 2.4 cm. Calculate the pipe's weight in lbs. (1 lb = 4.4482 N)

\textbf{Material}
\textbf{Density (kg/m$^{3}$)}

1
Copper
8.9 × 10$^{3}$

2
Iron
7.8 × 10$^{3}$

3
Lead
11.3 × 10$^{3}$



\emph{w }= \rule{2cm}{0.4pt} lbs

\item \hfill \textit{1 pts}

A cheerleader of mass 48 kg stands on the shoulders of a football player of mass 86 kg. The football player is standing in a soft, thin layer of mud that does not permit air under his shoes. If each of his shoes has an area of 279 cm$^{2}$, calculate the absolute pressure exerted on the surface underneath \textbf{\emph{one}} of his shoes. Assume \emph{g} = 9.80 m/s$^{2}$ and atmospheric pressure is 101,000 Pa.



\emph{p} =\rule{2cm}{0.4pt}Pa

\item \hfill \textit{1 pts}

Reading the scale on the mercury barometer in the figure, a researcher on Planet X finds that \emph{h}$_{1}$ = 6 cm and \emph{h}$_{2}$ = 51.1 cm. Calculate the ambient air pressure in kPa. Density of mercury: 13,534 kg/m$^{3}$. 1 kPa = 1000 Pa. The acceleration of gravity on Planet X is 2.6 m/s$^{2}$.

[Image: barometer01.png]


$p_0=\_\_\_\_\_\;\mathrm{kPa}$

\item \hfill \textit{1 pts}

Calculate the pressure at the bottom of a cylinder of mercury 29 cm deep if the ambient atmospheric pressure is 66,488 Pa. The density of mercury is 13,600 kg/m$^{3}$.



\emph{p} = \rule{2cm}{0.4pt} Pa

\item \hfill \textit{1 pts}

Suppose a ship's fuel tank has a layer of oil 1 m thick and density 617 kg/m$^{3}$ floating atop a layer of sea water of density 1,033 kg/m$^{3}$ that is 2.9 m thick. Given the ambient atmosphere has a pressure of 105 kPa, calculate the pressure at the bottom of the water layer in kPa.



\emph{p} = \rule{2cm}{0.4pt} kPa

\item \hfill \textit{1 pts}

Suppose a hydraulic press has an output piston that has 5.6 times the cross-sectional area as the input piston. If the input force is 314 N, what is the output force in kN?



\emph{F} = \rule{2cm}{0.4pt} kN

\item \hfill \textit{1 pts}

A small hydraulic press has input piston of radius 7 cm and output piston of radius 16 cm. The output piston is connected to a pressing plate of area 11.6 cm$^{2}$. If the input force is F$_{1}$ = 282 N, calculate the compressive stress in kPa exerted by the press.

[Image: HP4.png]

\emph{p} = \rule{2cm}{0.4pt} kPa

\end{enumerate}

\end{document}
